% ============================================================================
%  frontmatter.tex ------ 标题页、摘要、译者说明
%  本文件由主译者撰写，不含原文正文。
% ============================================================================

\begin{titlepage}
\centering
\vspace*{1.2cm}

{\LARGE\bfseries 费米路径积分与拓扑相\par}
\vspace{0.5cm}
{\large Fermion Path Integrals And Topological Phases\par}
\vspace{1.6cm}
{\large Edward Witten\par}
\vspace{0.35cm}
{\small School of Natural Sciences, Institute for Advanced Study,\\
Princeton, NJ 08540, USA\par}
\vspace{1.8cm}

{\large 中文译注版\par}
\vspace{0.4cm}
{\small 原文：arXiv:1508.04715v2 [cond-mat.mes-hall]，2015 年 12 月 19 日\par}
{\small 期刊版：Rev. Mod. Phys. \textbf{88}, 035001 (2016)\par}
\vspace{1.5cm}

{\small 全文逐段翻译，公式、编号与文献引用与原文一一对应\par}
\vfill
\end{titlepage}

% ----------------------------------------------------------------------------
%  摘要
% ----------------------------------------------------------------------------
\section*{摘要}
\addcontentsline{toc}{section}{摘要}

对称保护拓扑（symmetry-protected topological，SPT）物相已被用反常（anomaly）的语言加以刻画；人们也一直期望，对含费米子的 SPT 相应当存在一幅类似的图像。本文详细阐明：对于可由能带理论与自由费米子理解的量子物相，这幅图像究竟意味着什么。我们考虑的主要例子是 2 维或 3 维空间中的时间反演不变拓扑绝缘体与拓扑超导体。在此过程中，我们澄清了“三维拓扑绝缘体体内电磁 $\theta$ 角等于 $\pi$”这一陈述的确切含义。

% ----------------------------------------------------------------------------
%  译者说明
% ----------------------------------------------------------------------------
\section*{译者说明}
\addcontentsline{toc}{section}{译者说明}

\begin{enumerate}[label=\arabic*.,leftmargin=2.2em,itemsep=0.25em]
  \item \textbf{底本与版本}：译文以 arXiv:1508.04715v2（2015 年 12 月 19 日修订版）为底本。该版本共 89 页（含附录与参考文献），正文第 1--5 节、附录 A--C 全部译出；参考文献表按学术惯例保留英文原貌。
  \item \textbf{公式与编号}：所有公式均为原稿 LaTeX 源码的逐字移植，未作任何改写；交叉引用（\verb|\label|/\verb|\ref|）、文献引用（\verb|\cite|）、公式编号均沿用原稿，因此本书公式号与原文逐号对应，便于对照核查。若原文确有笔误或排印疏漏，一律保留原式，并在其后以 \verb|% [Tutor Note: ...]| 批注指出。
  \item \textbf{术语}：采用中文学界约定俗成的译法（如 gauge theory 规范场论、holonomy 和乐群、anomaly 反常、symmetry-protected topological phase 对称保护拓扑相、spectral flow 谱流、index theorem 指标定理、cobordism 配边、pin structure pin 结构等）。重要术语首次出现时以“中文（English）”形式给出，其后径用中文。
  \item \textbf{图表}：原文图 1--9 均自原 PDF 逐张裁出（\texttt{figs/fig01.png}--\texttt{fig09.png}），插入到对应段落；图注译为中文，并标注该图在原 PDF 中的页码与图号。
  \item \textbf{批注}：以 \verb|% [Tutor Note: ...]| 形式写在 LaTeX 注释中（编译后不显示，阅读源码可见），内容包括：易错的指标与张量缩并、作者省略的证明步骤、术语歧义与物理图像澄清。
  \item \textbf{本节重点}：每一大节（Section）开头以蓝框“本节重点（Key Takeaways）”给出该节的物理图像、数学工具与核心结论，系译者所加，非原文内容。
\end{enumerate}

\vspace{0.5em}
\noindent\textit{（下接原文正文。原文的目录页码与本版页码不同，正文结构一一对应。）}
